\documentclass[12pt]{article}

\usepackage[spanish]{babel}
\usepackage[utf8]{inputenc}
\usepackage[T1]{fontenc}
\usepackage{geometry}
\usepackage{amsmath, amssymb}
\usepackage{array}

\geometry{margin=2.5cm}
\begin{document}
\begin{titlepage}
\centering

\vspace*{\fill}
{\Large \textbf{Proyecto CA0305}} \\
{\large Bitácora 1} \\
\vspace{2cm}
{\large \textbf{Integrantes}} \\
\vspace{1cm}
\begin{tabular}{p{6cm} p{2cm}}
Ashly Garro Villanueva & C33198 \\
Robert Meléndez Jiménez & C34772 \\
Adriana Monge Rodríguez & C35042 \\
Emily Sánchez Mancía & C27260 \\
Alessandro Umaña Vega & C37963 \\
\end{tabular}

\vspace{2cm}

{\large \textbf{I Ciclo 2026}}
\vspace*{\fill}

\end{titlepage}

\newpage

% -------------------------
\section*{Fundamentación Matemática}

\subsection*{Task 01: Tensor Broadcasting Basics}

\subsection*{Task 02: Matrix Multiplication (Naive vs Vectorized)}
Las operaciones entre matrices son permite combinar dos matrices y obtener una nueva, se tiene dos matrices A
 y B , donde las dimensiones de A son (N, K) y las de B son (K, M), entonces definimos la nueva
  matriz C como c=AB, note que es importante verificar que las columnas de A y las filas de B sean del mismo tamaño
  para poder realizar la multiplicación. Dado la multiolicación las dimensiones de la matriz C serían (N,M).\\
Para poder calcular las entradas de la matriz C lo que se realiza es que cada entrada(i,j) de la matriz se calcula
 mediante el producto punto, donde se multiplica la fila A por la columna B y se suma cada entrada de esa multiplicación.
 De esta manera se consigue cada entrada de la matriz C. Tenemos dos formas de realizar este cálculo, mediante
  la fórmula matemática con la sumatoria de la multiplicación de las entradas y usando las funciones numpy
   que hacen multiplicación vectorial de forma automática.
\subsection*{Task 03: Element-wise Operations}

\subsection*{Task 04: Tensor Reshaping \& Transposing}

\subsection*{Task 05: Reduction Operations}
Las operaciones de reducción son funciones que toman un tensor de gran dimensión
y lo transforman en uno de menor dimensión, colapsando una o más dimensiones mediante alguna operación. \\
Un tensor es un arreglo multidimensional de números, un ejemplo de esto es una matriz de dos dimensiones. 
La idea central al colapsar estos tensores es que, en muchas ocasiones, se tienen una gran cantidad de valores
organizados en varias dimensiones, pero se necesita resumirlos en algo más simple y manejable. \\
Las cuatro operaciones implementadas corresponden a distintas maneras de colapasar estos tensores. La suma acumula
todos los elementos de la dimensión especificada. La media divide esa suma entre la cantidad de elementos de esa dimensión, lo que la hace
independiente del tamaño del arreglo. El máximo selecciona el valor más grande de esa dimensión. Finalmente, el argmax está relacionado
con el máximo pero en lugar de devolver el valor más grande, devuelve el índice donde se encuentra ese valor. \\

\subsection*{Task 06: Vector Norms (L1, L2)}
Una norma es una función que toma un vector y cuantifica su tamaño o "longitud". En este caso, la longitud no se 
refiere a la longitud física, sino a una medida de qué tan grande es el vector en términos de sus componentes. \\
La norma L1 mide el tamaño de un vector sumando los valores absolutos de todos sus componentes: $\displaystyle \sum_{i=1}^{D} \mid x_i \mid$.
El valor absoluto es necesario ya que sin él los valores negativos cancelarían los positivos y el resultado no representaría una magnitud real.\\
En la norma L2 \Bigg($\displaystyle \sqrt{\sum_{i=1}^{D} x_i^2}$\Bigg), el cuadrado elimina los signos negativos igual que el valor absoluto en L1, y la raíz cuadrada al final devuelve el resultado a las
unidades originales. Geométricamente mide la distancia en línea recta desde el origen hasta el punto que representa el vector, que es exactamente la 
distancia euclideana clásica. \\
Para entender la importancia de estas normas en la práctica, es importante entender qué son los pesos y cómo se entrenan. Los pesos son números internos
de la red que determinan cómo procesa la información. Al principio se inicializan con valores aleatorios y la red no sabe nada. El entrenamiento es el 
proceso de ajustar esos pesos poco a poco, la red procesa una entrada, se calcula el error que cometió, se determina en qué dirección ajustar cada peso 
para reducir ese error, y se aplica ese ajuste. Luego de repetir este ciclo millones de veces con muchos ejemplos, los pesos van tomando valores que hacen
que la red produzca respuestas cada vez más correctas.\\
El uso de las normas empieza a lo largo de este proceso, durante este, los pesos puede crecer en gran magnitud, llevando a que la red memorice los datos en vez
de aprender patrones. Las normas permiten medir y controlar ese crecimiento de dos maneras. La primera es la regularización, donde al calcular el loss (qué tan equivocada
estuvo la red), se agrega un término extra basado en la norma de los pesos que castiga valores grandes (si los valores son grandes el loss aumenta y la red no lo pudo minimizar),
obligando a la red a mantener pesos pequeños mientras minimiza el error. \\
La segunda es el gradient clipping, donde se calcula la norma L2 de los ajustes que se le van a aplicar a los pesos y si se supera un umbral se reducen proporcionalmente antes de
 aplicarlos, evitando que el entrenamiento se desestabilice por actualizaciones demasiado grandes. 
\subsection*{Task 07: Dot Product \& Cross Product}

\subsection*{Task 08: Einstein Summation (einsum) Basics}

\subsection*{Task 09: Advanced Einsum (Batch MatMul)}

\subsection*{Task 10: Gradient of Sum}

\subsection*{Task 11: Gradient of Matrix Multiplication}

\subsection*{Task 12: One-Hot Encoding}

\subsection*{Task 13: Softmax Implementation (Stable)}
La función softmax es una transformación que toma un vector de valores reales
arbitrarios (llamados logits) y los convierte en una distribución de probabilidad 
válida, lo que permite interpretarlos como la probabilidad de que una muestra 
pertenezca a cada clase.\\
Su fórmula es $\displaystyle \frac{e^{x_i}}{\sum_{j=i}^{C} e^{x_j}}$ donde $x_i$ es el 
logit de la clase $i$ y $C$ es el número total de clases. En primer lugar, la exponencial convierte cualquier
valor real en un número positivo, lo que garantiza probabilidades positivas. 
Segundo, aumenta las diferencias entre logits de forma no lineal, si un logit es ligeramente mayor que los demás,
su exponencial será más grande en gran magnitud, produciendo una probabilidad dominante. Al dividir entre la suma
de todas las exponenciales, se normaliza el resultado para que todo sume 1.\\
En el contexto de redes neuronales, el softmax ocupa un lugar importante en la estructura de cualquier red de clasificación.
Todo lo que hacen las tareas anteriores (procesar entradas, extraer características, aprender patrones), termina produciendo
un vector de logits que por sí solos no son interpretables. El softmax es la última capa que convierte esos logits en probabilidades
con significado real e interpretable. Sin este, o sería posible calcular el cross-entropy loss, que es el número que mide el error de
la red y guía el ajuste de sus pesos durante el entrenamiento. En ese sentido el softmax es el puente entre las operaciones internas de 
la red y la salida interpretable.
\subsection*{Task 14: Cross-Entropy Loss (Manual)}
La cross-entropy es una función de perdida utilizada principalmente para cuantificar que tan diferentes son dos distribuciones que predicen un mismo modelo.
Este modelo se enfoca en la probabilidad asignada a la clase correcta, su función principal es que califica el modelo según cuanta probabilidad le asigna a la respuesta correcta, si se le asigna mucha probabilidad entonces el error es pequeño por otro lado si se le asigna muy poca probabilidad el error crece mucho. En vez de solo decir si el resultado es correcto o incorrecto este modelo mide que tan acertado o equivocado esta el modelo simultáneamente, el objetivo de este modelo es que vaya aumentado la probabilidad de respuesta correcta conforme aumentan los intentos.
\subsection*{Task 15: Numerical Stability (Log-Sum-Exp)}
El truco de log-sum Exp es una técnica para poder calcular expresiones de la forma \[
\log\left(\sum_{j=1}^{D} e^{x_j}\right)
\]
De forma que cuando se evalúan números muy grandes estos no se vayan hacia infinito y se pueda llegar a una aproximación realista. Si se calculara sin este truco la gran mayoría de veces daría infinito, para esto se usa la identidad \[
\log\left(\sum_{j=1}^{D} e^{x_j}\right)
=
m+\log\left(\sum_{j=1}^{D} e^{x_j-m}\right)
\]
La idea es que la suma de exponenciales puede reescribirse factorizando el valor máximo de los exponentes. Al hacer esto, se extrae ese máximo fuera de la suma y se convierte en un término separado. Luego, usando la propiedad del logaritmo de que el logaritmo de un producto es igual a la suma de los logaritmos, se separa la expresión en dos partes: el máximo y el logaritmo de una suma de exponentes ajustados.
Donde denotamos m como el max de los exponentes de todas las exponenciales, así si restamos este a todas las otras, su mayor valor puede ser 1 aprovechando las propiedades que tenemos sobre $e$ cuando tiende a -infinito. Se utiliza N como cantidad de filas y D cantidad de valores por fila, el procedimiento para realizar este truco es encontrar el max de cada fila, restarlo para ajustar los exponentes de la e, aplicar exponencial y sumar por fila, aplicamos el logaritmo y sumamos nuevamente el máximo encontrado. El resultado final es un vector de tamaño (N, ) con el log-sum-exp de cada fila.

% -------------------------
\section*{Conclusiones}

\end{document}